\clearpage\SourceLeaf{231}\ArithmeticHead{146}{L}{L. SCHLESINGER, ÜBER GAUSS' FUNKTIONENTHEORETISCHE ARBEITEN.}
\indent Der erste Abschnitt der Abhandlung (Artikel 7—11) gibt die linearen
Beziehungen mit von \(x\) abhängenden Koeffizienten, die zwischen je drei Reihen
\[
\begin{gathered}
F(\alpha,\beta,\gamma,x);\quad F(\alpha+\lambda,\beta+\mu,\gamma+\nu,x);\\
F(\alpha+\lambda',\beta+\mu',\gamma+\nu',x)
\end{gathered}
\]
bestehen, wenn \(\lambda,\mu,\nu,\ \lambda',\mu',\nu'\) die Werte \(0,+1,-1\) haben (Art. 11). Aus
diesen Beziehungen zwischen »verwandten Funktionen« (vergl. besonders die
Form (11) oben S. 139) ist der für die neuere Theorie der linearen Differentialgleichungen
so bedeutungsvolle \textls[100]{\textsc{Riemannsche} Klassenbegriff} hervorgegangen.
\par
Im zweiten Abschnitt (Artikel 12—14) werden die Kettenbruchentwicklungen
für den Quotienten (9) mit Anwendung des auf \textsc{Lambert} zurückgehenden
Divisionsverfahrens hergeleitet, ohne dass jedoch in eine Erörterung der Konvergenzfrage
eingetreten wird. Hier werden die in der \textit{Theoria motus}, Art. 90
gegebenen Formeln eingeordnet.
\par
Der dritte Abschnitt bringt (Artikel 15—18) die Konvergenzuntersuchung\linebreak
von \(F(\alpha,\beta,\gamma,1)\), die auch durch Reihenvergleichung, nämlich mit der Reihe
\[
1+\frac{n-h-1}{n}+\frac{(n-h-1)(n-h)}{n(n+1)}+\ldots=\frac{n-1}{h},
\]
(Werke III, S. 142) geführt wird. \textsc{Gauss} beschränkt sich dabei auf reelle
Werte der \(\alpha,\beta,\gamma\), seine Methode lässt aber eine unmittelbare Übertragung
auf den Fall komplexer Werte zu\footnote{Siehe K. \textsc{Weierstrass}, \textsc{Crelles} Journal 51 (1856), S. 22ff.}. Um den Wert von \(F(\alpha,\beta,\gamma,1)\) als
Funktion der \(\alpha,\beta,\gamma\) darstellen zu können, wird dann (Artikel 18—20) auf
\textsc{Eulersche} Weise (vergl. oben) die Funktion \(\Pi(z)\) eingeführt, wobei die exakte
Konvergenzuntersuchung des Produkts (10) besonders erwähnt werden mag.
Da \(\Pi(z)\) für jedes reelle \(z\) bekannt ist, wenn man seine Werte für \(0<z<1\)
kennt (Art. 23), wird eine Tafel der letzteren Werte auf 20 Dezimalstellen
für die Hundertteile des Intervalls \(0\ldots1\) gegeben. Es folgt (Art. 24) die
Darstellung von \(F(\alpha,\beta,\gamma,1)\) durch die \(\Pi\)-Funktion, (Artikel 25, 26) die beiden
Funktionalgleichungen für \(\Pi(z)\), (Art. 27) das sogenannte \textsc{Eulersche} Integral
erster Gattung (die \textsc{Legendresche} \(B\)-Funktion) mit Anwendung auf die Berechnung
der lemniskatischen Periode\footnote{Vergl. die Aufzeichnung in der Scheda Aa (1798), abgedruckt Werke III, S. 413 und die Aufzeichnung auf dem Deckel des \textit{Leiste}, Werke X, 1, S. 145.}, (Art. 28) das \textsc{Eulersche} Integral
zweiter Gattung, (Art. 29) die semikonvergente Reihe für \(\log\Pi(z)\)\footnote{In bezug auf diese Art von Reihen heisst es Werke III, S. 152: »Im übrigen kann man nicht leugnen, dass die Theorie dieser Art von divergenten Reihen bisher mit gewissen Schwierigkeiten einhergeht, über die wir vielleicht bei anderer Gelegenheit ausführlich handeln werden.«}, (Artikel\par
\clearpage\SourceLeaf{232}\ArithmeticHead{147}{R}{VIIc. ANALYSE DER DISQUISITIONES.}
\noindent 30, 31) die Einführung der Funktion
\[
\Psi(z)=\frac{d\log\Pi(z)}{dz}
\]
und Berechnung der \textsc{Euler-Mascheronischen} Konstanten \(-\Psi(0)\) durch die
semikonvergente Entwicklung, (Artikel 32, 33) einige Anwendungen der Funktion
\(\Psi\), (Artikel 34, 35) ihre Integraldarstellung, endlich (Artikel 36, 37) eine
Konvergenzuntersuchung der für \(\Pi,\ \Pi\Psi\) und \(\Psi\) gefundenen Integrale. In bezug
auf diese bemerkt \textsc{Gauss} (Werke III, S. 159): »Dieses Beispiel zeigt,
welche Vorsicht erforderlich ist, wenn man es mit unendlichen Grössen zu
tun hat; nach unserer Überzeugung dürfen solche bei mathematischen Untersuchungen
nur in soweit zugelassen werden, als sie sich auf die Theorie der
Grenzwerte reduzieren lassen.«
\par
Überblicken wir diesen Inhalt der \textit{Disquisitiones}, so drängen sich zwei
Bemerkungen auf. Erstens vermissen wir die \textsc{Eulersche} Darstellung von
\(F(\alpha,\beta,\gamma,x)\) durch ein bestimmtes Integral für beliebige Werte der \(\alpha,\beta,\gamma\);
zweitens muss uns, die wir \textsc{Gauss} nicht nur nach seinen Publikationen, sondern
nach seinen handschriftlichen Aufzeichnungen kennen, die ganze Art
der Darstellung befremdlich erscheinen, die ganz im Stile des XVIII. Jahrhunderts
gehalten ist und mit ängstlicher Scheu allem aus dem Wege geht,
was eine prinzipielle Anwendung der komplexen Variabeln erfordern würde.
In dem Briefe an \textsc{Bessel} vom 18. Dezember 1811 (Werke X, 1, S. 366) lesen
wir: »Zuvörderst würde ich jemand, der eine neue Funktion in die Analyse
einführen will, um eine Erklärung bitten, ob er sie bloss auf reelle Grössen
\mbox{.\ .\ .}\ angewandt wissen will, und die imaginären Werte des Arguments gleichsam
nur als ein Überbein ansieht, oder ob er meinem Grundsatz beitrete,
dass man in dem Reiche der Grössen die imaginären \(a+bi\) als gleiche Rechte
mit den reellen geniessend ansehen müsse. Es ist hier nicht von praktischem
Nutzen die Rede, sondern die Analyse ist mir eine selbständige Wissenschaft,
die durch Zurücksetzung jener fingierten Grössen ausserordentlich an Schönheit
und Ründung verlieren und alle Augenblick Wahrheiten, die sonst allgemein
gelten, höchst lästige Beschränkungen beizufügen genötigt sein würde.«
\textsc{Gauss} ist also in den \textit{Disquisitiones} diesem seinem eigenen Grundsatz nicht\par
\HistoricalFoot{}{19\textsuperscript{*}}
\clearpage\SourceLeaf{233}\ArithmeticHead{148}{L}{L. SCHLESINGER, ÜBER GAUSS' FUNKTIONENTHEORETISCHE ARBEITEN.}
\noindent »beigetreten«, ja wir werden wohl nicht fehlgehn, wenn wir die Teilung der
Abhandlung gerade an dieser Stelle, mit dem Umstande in Verbindung bringen,
dass, wie wir sogleich sehen werden, der erste Artikel der von \textsc{Gauss} nicht
mehr veröffentlichten Fortsetzung wenigstens andeutungsweise mit Wegen der
Variabeln \(x\) im komplexen Gebiete operiert. Vielleicht wird sich eine Erklärung
dieser beiden Tatsachen ermöglichen lassen, wenn wir versuchen, uns
über den Umfang von \textsc{Gauss'} funktionentheoretischen Kenntnissen zu orientieren.
Wir nehmen zu dem Ende den Bericht über die Abhandlung, bezw.
deren nachgelassene Fortsetzung jetzt wieder auf.\par

